\cleardoublepage

\section{相关工作}

\subsection{基于大模型的单元测试生成}

基于大语言模型的单元测试生成通常将待测代码、函数签名、注释信息以及项目上下文组织为提示词输入模型，再由模型输出测试代码、断言语句或测试样例。与传统自动化测试生成方法相比，大模型方法具有较强的语义泛化能力，能够在较少人工规则约束的情况下生成更贴近自然语言描述的测试代码。

现有方法大致可以按照输入组织方式和生成策略进行划分。部分工作采用纯提示驱动方式，仅通过代码片段和自然语言说明生成测试；部分工作结合检索增强机制，从项目历史测试、接口文档或相似代码中检索上下文，再交由模型生成；还有工作引入反馈闭环，在测试编译或执行失败后将错误信息反馈给模型进行迭代修复。这些方法在测试可读性、覆盖率和生成效率方面取得了一定进展，但也面临输出不稳定、复杂依赖处理能力有限以及对模型调用成本较高等问题。

此外，单元测试生成通常不仅要求代码语法正确，还要满足测试框架、依赖导入、初始化顺序和断言逻辑等约束。对于复杂项目而言，模型容易出现遗漏依赖、断言不准确或测试目标偏移的问题。因此，如何在生成阶段结合结构化程序信息，并在生成后引入自动化校验，是提升单元测试质量的重要方向。

\subsection{基于大模型的代码修复}

代码修复是大语言模型在软件工程场景中的另一项重要应用。针对编译错误、静态分析告警和运行时异常，研究者通常会将错误信息、堆栈信息、相关代码片段与补充上下文一并输入模型，生成修复补丁或重写后的代码版本。与人工修复相比，大模型方法能够快速给出候选修复方案，特别适合用于测试代码生成后的自动纠错阶段。

现有代码修复方法主要包括直接补丁生成、基于错误定位的局部重写以及多轮交互式修复。直接补丁生成实现简单，但容易引入过度修改；局部重写能更好地控制修改范围，但依赖准确的错误定位；多轮交互式修复则能够在反馈驱动下逐步提升修复结果，但会引入额外调用开销。对于单元测试场景，修复任务不仅要保证语法和类型正确，还要兼顾测试意图、断言语义以及可重复执行性，因此仍然需要更细粒度的错误分析与模型协同策略。

综合来看，现有基于大模型的单元测试生成与代码修复研究已经证明了语言模型在软件测试任务中的潜力，但在“生成”和“修复”两个关键阶段仍然存在资源利用率不足、模型分工不清和性能提升受限等问题。本文提出的 HieraTest 正是围绕这些不足展开设计。
